%
% 2-verallgemeinert.tex
%
% (c) 2026 Prof Dr Andreas Müller
%
\section{Verallgemeinerte Potenzreihen}
\kopfrechts{Verallgemeinerte Potenzreihen}
Eine Potenzreihe hat innerhalb des Konvergenzradius keine Singularitäten.
Andererseits ist es durchaus möglich, dass eine Differentialgleichung
auf einem Intervall, welches $0$ enthält, eine Lösung hat, die sich
nahe an $0$ wie eine Wurzelfunktion oder wie eine Potenz von $1/x$
verhält.
Die Ableitung solcher Funktionen wächst für $x\to 0$ über alle
Grenzen.
Eine Potenzreihe hat dagegen immer beschränkte Ableitung in einer 
genügend kleinen Umgebung des Entwicklungspunktes.
Um eine grössere Klasse von Differentialgleichung mit einer Reihe
lösen zu können, muss daher die Klasse der zulässigen Reihen etwas
vergrössert werden.

\begin{definition}[verallgemeinerte Potenzreihe]
\label{buch:differentialgleichungen:verallgemeinert:def}
Eine \emph{verallgemeinerte Potenzreihe} ist ein Reihe der Form
\index{verallgemeinerte Potenzreihe}%
\begin{equation}
f(z)
=
\sum_{k=0}^\infty
a_k z^{k+\varrho}
\qquad\text{oder}\qquad
f(z)
=
\sum_{k=0}^\infty
a_k (z-z_0)^{k+\varrho}
\label{buch:differentialgleichungen:verallgemeinert:eqn:def}
\end{equation}
mit $a_0\ne 0$ und $\varrho\in\mathbb{R}$.
\end{definition}

Eine verallgemeinerte Potenzreihe mit $\varrho=0$ ist eine gewöhnliche
Potenzreihe.
Für $0<\varrho<1$ entsteht eine im Punkt $0$ (bzw.~$z_0$)
stetige Funktion, deren Ableitung für $z\to 0$ (bzw.~$z\to z_0$)
über alle Grenzen wächst.


Eine verallgemeinerte Potenzreihe $f(z)$ kann auch als
\begin{equation}
f(z)
=
z^{\varrho}
\sum_{k=0}^\infty
a_k z^k
\label{buch:differentialgleichungen:verallgemeinert:eqn:gewoehnlich}
\end{equation}
geschrieben werden.
Die Reihe 
\eqref{buch:differentialgleichungen:verallgemeinert:eqn:def}
ist daher genau dann konvergent, wenn die gewöhnliche Potenzreihe
auf der rechten Seite von
\eqref{buch:differentialgleichungen:verallgemeinert:eqn:gewoehnlich}
konvergiert.
Der Konvergenzradius ändert sich beim Übergang zu einer verallgemeinerten
Potenzreihe also nicht.

Die $l$-te Ableitung einer verallgemeinerten Potenzreihe ist
\begin{align*}
f^{(l)}(z)
&=
\sum_{k=0}^\infty
(k+\varrho)
(k+\varrho-1)
\cdots
(k+\varrho-l+1)
a_k z^{k+\varrho-l}
\\
&=
z^{\varrho-l}
\sum_{k=0}^\infty
\underbrace{
(k+\varrho)
(k+\varrho-1)
\cdots
(k+\varrho-l+1)
a_k}_{\displaystyle=b_k} z^k.
\end{align*}
Dies ist wieder eine verallgemeinerte Potenzreihe mit dem neuen
Exponenten $\varrho-l$.
Der Konvergenzradius kann mit der Formel
\begin{align*}
\frac{1}{r_{f^{(l)}}}
&=
\limsup_{n\to\infty}
\!\sqrt[n]{
|(n+\varrho)(n+\varrho-1)\cdots(n+\varrho-l+1)|
\cdot
|a_n|
}
\\
&\le
\limsup_{n\to\infty}
\!\sqrt[n]{
\raisebox{1pt}{$\mathstrut$}
\smash{
(n+\varrho)^l
\cdot
|a_n|
}
}
=
\lim_{n\to\infty}
(n+\varrho)^{\frac{l}{n}}
\cdot
\limsup_{n\to\infty}
\!\sqrt[n]{
|a_n|
}
\\
&=
\lim_{n\to\infty}
n^{\frac{l}{n}}
\cdot
\frac{1}{r_{f}}
=
\frac{1}{r_f}
\end{align*}
von Cauchy-Hadamard bestimmt werden.
Die Ungleichung ist erfüllt, sofern $n$ viel grösser als $|\varrho|+l$
ist.
Der Konvergenzradius $r_{f^{(l)}}$ der $l$-ten Ableitung $f^{(l)}(z)$
ist also gegenüber dem Konvergenzradius $r_f$ der ursprünglichen Funktion
$f(z)$ unverändert.


